\documentclass[11pt]{article}
\usepackage[utf8]{inputenc}
\usepackage[T1]{fontenc}
\usepackage[margin=1in]{geometry}
\usepackage{amsmath,amssymb}
\usepackage{booktabs}
\usepackage{xcolor}
\usepackage[colorlinks=true,linkcolor=blue,urlcolor=blue]{hyperref}
\newcommand{\tr}{\operatorname{tr}}
\title{\textbf{Blind Recovery for Objective-Aware Scheduling}\\[2pt]
\large A bridge note: what GO-EC-3's 94.6\% means if you don't care
about water-filling}
\author{Observation Theory program --- systems bridge note}
\date{August 2026 --- companion to ledger row \texttt{GO-EC-3}
(\texttt{GO-P-2026-087}, sealed, ALL PASS 5/5).}
\begin{document}
\maketitle

\section*{The problem, in systems terms}

You run a state estimator --- a Kalman filter, say --- over a pool of
candidate sensors, and each step you can only afford to sample a few of
them. Which few? The textbook answers maximize information: pick the
sensors that shrink the estimate covariance the most (\emph{logdet}), or
that cut its trace (\emph{iso-trace}). Both are \textbf{consumer-blind}:
they optimize the estimate in the abstract, not the thing that actually
\emph{reads} the estimate downstream --- your controller, your detector,
your serving policy, whatever consumes the state and has an objective of
its own. When the downstream objective is anisotropic (it cares about
some directions of the state far more than others), consumer-blind
scheduling spends your sensor budget in the wrong place.

The catch: you usually can't \emph{open} the downstream consumer. It may
be a learned model, a third-party component, or a policy too complex to
differentiate by hand. So the question is: \textbf{can you schedule
sensors for a downstream objective you can only probe as a black box ---
and is it worth the probing cost?}

\section*{The idea (no water-filling required)}

Treat the consumer as a black box you can \emph{query}: feed it a state
estimate, read back its loss. With a handful of finite-difference queries
you recover a local model of what it cares about --- a positive
semidefinite \emph{read operator} $\hat P_C$, the curvature of its loss
in state space. Then you score any candidate sensor set by how much it
reduces the error the consumer will feel, which is the single scalar
\[
  \tr\!\big(\hat P_C\,\Sigma\big),
\]
$\Sigma$ being the Kalman covariance you already track. Greedily pick the
sensors that drive this down. That's the whole pipeline: probe, compose,
schedule. No knowledge of the consumer's internals; no analytic
objective; and --- crucially --- every probe query is \textbf{charged
against the same sensor budget}, so the comparison is honest.

\section*{What GO-EC-3 measured}

A single sealed, pre-registered governed run (fresh seed \texttt{20260819},
20 fresh systems, probe cost charged), two consumer arms:

\begin{itemize}
\item \textbf{Positive control (analytic consumer).} Here the optimal
schedule is \emph{known} in closed form, so the blind pipeline has a
target to match, not beat. It captured \textbf{94.6\%} of the known
optimum's gain (gate $\ge 75\%$). In plain terms: probing recovered
almost all of what you would get if you had the consumer's source code.
\item \textbf{Black-box consumers (no analytic optimum).} A frozen
low-rank MLP and a smoothed-threshold consumer, where no closed form
exists. The blind schedule beat the best consumer-\emph{agnostic} policy
by \textbf{16.3\%} at matched budget (gate $\ge 8\%$).
\item \textbf{Ordering.} Across 61 system pairs, the blind score
predicted which schedule wins \textbf{86.9\%} of the time (gate
$\ge 65\%$).
\item \textbf{No free lunch.} Probing a consumer with \emph{shuffled}
inputs gave only $+0.08$ (a near-null control); an \emph{anti} policy
that schedules against the recovered operator was worst per arm; the
probe charge itself cost $+0.0030$ (a sensitivity row, no verdict
weight).
\end{itemize}

\begin{center}
\small
\begin{tabular}{llll}
\toprule
gate & quantity & result & bar \\
\midrule
P1 & match of known optimum (positive control) & 0.946 & $\ge 0.75$ \\
P2 & improvement over best agnostic (black box) & 0.163 & $\ge 0.08$ \\
P3 & trace-matched ordering, 61 pairs & 0.869 & $\ge 0.65$ \\
\multicolumn{2}{l}{shuffled-consumer / anti controls} & clean & --- \\
\multicolumn{2}{l}{probe charge (sensitivity)} & $+0.0030$ & --- \\
\bottomrule
\end{tabular}
\end{center}

\section*{Why a systems person should care}

The marketable claim is not about geometry. It is: \textbf{you can make a
scarce sampling budget serve a downstream objective you cannot open,
automatically, by querying it --- and recover almost all of the gain you
would get from white-box access, while paying for the probes.} The same
recipe --- probe the consumer, compose with what you already track,
schedule --- is objective-agnostic in the useful sense: swap the Kalman
covariance for a cache's access statistics and you have objective-aware
eviction; swap it for a feature store's uncertainty and you have
query scheduling. It turns ``we don't know what the downstream wants''
from a blocker into a bounded probing cost.

\section*{What it is not (scope, stated plainly)}

The systems in GO-EC-3 are \textbf{synthetic}: linear-time-invariant
dynamics with planted consumers. The greedy scheduler carries \textbf{no
optimality claim} (submodular greedy is a heuristic here). It is a
\textbf{single} governed run, reported per its pre-registration. The
result earns the \emph{mechanism} --- blind recovery works, and it pays
for itself --- not a deployment number. The physical-endpoint version,
on a real estimator and a real downstream, is a separate campaign
(Campaign~5, its own registry ID) and has not been run. Read this note as
``here is why the mechanism is worth a real trial,'' not ``here is a
benchmark.''

\end{document}
